% ===========================================================================
% Section 4: the completeness result.
% Sources: paper_rewrite_plan.md §2 thesis item 2, §6 subplan "§3 Why Basis
% Identity Cannot Matter -- and Where It Does", Part II §20 (evidence
% inventory rows "Basis identity at complete rank" and "Nonlinearity
% placement"); rebuttal responses response_K8yb.md (i), response_CAcc.md §2,
% response_ovPT.md Q4, response_y8d1.md W2, response_AC.md opening + Point 4.
%
% Statistics quoted exactly: rotation-absorption residual 5.6e-15; 36 paired
% comparisons, none significant, minimum p = 0.084; boundary test KW
% p = 6.4e-05 (fully linear), p = 0.457 (after basis), p = 1.9e-07
% (interposed). The four scope conditions restate the enumerate of
% Sec.~method_heads verbatim in substance.
% ===========================================================================

\section{Basis Completeness: Why Basis Identity Cannot Matter}
\label{sec:completeness}

The trunk result removes width from the story; this section removes the
basis identity from it, at complete forecast rank. What remains when both
are removed --- rank and captured energy --- is the subject of
Section~\ref{sec:predictors}.

\subsection{The rotation-absorption argument}
\label{sec:completeness:proof}

The argument is a two-line consequence of the formulation
\eqref{eq:formulation}. Suppose $d_f = H$, the paper's non-tiered
default. Then $\mathbf{B}_f \in \mathbb{R}^{H \times H}$ is square and
orthogonal, and for any unconstrained synthesis head
$\mathbf{M} \in \mathbb{R}^{H \times U}$ there exists a learned projection
$\mathbf{W}_f = \mathbf{B}_f \mathbf{M}$ such that
$\mathbf{B}_f^{\top} \mathbf{W}_f = \mathbf{B}_f^{\top}
\mathbf{B}_f \mathbf{M} = \mathbf{M}$ exactly. A block with a complete
frozen forecast basis therefore represents precisely what standard N-BEATS
with an unconstrained head represents, in rotated coordinates: the freely
learned projection absorbs the rotation between any two complete
orthonormal bases, and on the forecast path basis identity cannot change
what a block can synthesize. We verified this numerically at residual
$5.6\times10^{-15}$.

Three null results that the submitted version reported without explaining
follow immediately: the family null, the filter-length null, and blocks
being interchangeable at equal rank. The limiting case is instructive and,
as Section~\ref{sec:method_basis} discloses, was actually shipped: below
its minimum filter length a wavelet generator returns the identity matrix,
the degenerate basis for which the absorption holds literally.

\subsection{Scope conditions}
\label{sec:completeness:scope}

Four scope conditions keep the claim honest; they are stated with the
formulation in Section~\ref{sec:method_heads} and every downstream use in
this paper respects them.
\begin{enumerate}
\item \textbf{Expressiveness, not trainability.} The guarantee fixes what
the forecast head can represent, not what training finds; a complete
basis whose coefficient spectrum is degenerate can fail to train on every
seed (Section~\ref{sec:predictors}).
\item \textbf{Trend-branch-present null.} Every arm of the identity
comparison grid carried the paired polynomial trend branch, which absorbs
the dominant component that otherwise breaks coefficient conditioning;
the measured null below is a trend-branch-present result.
\item \textbf{Forecast head only.} Every shipped backcast basis is
truncated ($d_b = H$ or $2H$ rows against a backcast of three to five
horizons), and truncated bases from different generators span different
subspaces that no rotation relates; on the backcast path family
equivalence rests on measurement, not proof.
\item \textbf{Tiered configurations excluded.} Tiered offset schedules
truncate the forecast basis itself and belong to the truncated regime of
Section~\ref{sec:predictors}, where measurement carries whatever
equivalence holds.
\end{enumerate}
\subsection{The measurement}
\label{sec:completeness:measurement}

The proposition is testable, and it was tested against the comparison a
reviewer named directly: Daubechies-3 \citep{daubechies1992ten} against a
frozen DCT and two random orthonormal draws, at equal rank, with
architecture, data, protocol, seeds, and parameter counts identical ---
byte-identical counts --- across all six M4 periods at ten seeds. The
result is 36 paired comparisons with none significant; the smallest raw
$p$ is 0.084. By construction every arm carried the paired trend branch,
so by scope condition~2 this is a trend-branch-present result. By the
reviewer's own criterion the contribution of the submitted version is
frozen low-rank projection rather than wavelets.

One disclosure sharpens rather than weakens the result. Because several
historical ``wavelet'' configurations silently received the identity basis
on their forecast path (coif2 and sym10 on five of six M4 periods,
Daubechies-3 on Yearly and Quarterly; only haar and db2 were genuine
throughout), not one of the six per-period M4 leaders carries a genuine
wavelet on its forecast path. The accuracy stands and the labels do not.
This weakens the family null as evidence about wavelets and strengthens
the proposition: the identity is its limiting case, and those cells score
exactly where the genuine wavelets beside them score. Accordingly the
family-equivalence analysis is restated over the families genuinely
distinct on each period (Hourly: all five genuine; Yearly and Quarterly:
\{haar, db2, identity\}; Weekly, Daily, Monthly: \{haar, db2, db3,
identity\}), every affected table cell carries the $\ddagger$ marker of
the unified legend scoped to exactly those (family, horizon) combinations,
and no sentence attributes an effect to a named family on an
identity-marked cell. The full disclosure is carried in
Table~\ref{tab:corrections_ledger}.

\subsection{Boundary test: where the premise is needed}
\label{sec:completeness:boundary}

The argument's premise is that nothing nonlinear sits between the learned
projection and the frozen basis. A 240-run placement study localizes
exactly where it is needed. With fully linear synthesis the Kruskal--Wallis
test rejects at $p = 6.4\times10^{-5}$, but the entire spread is the DCT
conditioning pathology of Section~\ref{sec:predictors} --- a trainability
failure, not an expressiveness difference. Placing the nonlinearity after
the basis restores the null ($p = 0.457$, spread 0.133). Interposing it
between projection and basis breaks the equivalence violently
($p = 1.9\times10^{-7}$, spread 40.487). The guarantee holds precisely
when its hypothesis does.

\subsection{Why the wavelet premise was misconceived}
\label{sec:completeness:misconceived}

Which property of compactly supported orthogonal dyadic wavelets
\citep{mallat1989wavelet} does this architecture exploit? None, and the
reason is a mismatch of approximation setting. The classical superiority
of wavelets is a theorem about nonlinear best-$m$-term approximation, in
which the retained coefficients are chosen per signal. Here the rows are
fixed at build time and shared across every series: linear approximation,
where the optimal fixed subspace is the Karhunen--Lo\`eve basis and where
localisation --- the property wavelets are chosen for --- becomes a
liability, because each retained translate covers only part of the time
axis. This answers the positioning question for short-time Fourier
transforms, wavelet packets, stationary transforms, and overcomplete
dictionaries in one stroke: the time-frequency versus time-scale axis is
the wrong axis for this architecture, and what governs the outcome is how
much target energy lies in the span of the retained rows, whatever
produced them (Section~\ref{sec:predictors}).

One corrected claim belongs to this paragraph as history. After
submission we claimed that a band-ordered wavelet finishes behind random
orthonormal draws; the claim was an artifact of our own
singular-value-decomposition orthogonalization, whose energy-descending
row ordering destroyed the dyadic structure. Rebuilt so that the band
ordering survives, the wavelet lands within 0.30 sMAPE of the DCT at
matched rank, and the corrected ordering lifts captured energy from 0.348
to 0.9999. The withdrawal is carried in
Table~\ref{tab:corrections_ledger}.

